\customchapter{INTRODUCCIÓN} 

\introsection{Presentación del tema de investigación}

\noindent El desarrollo de sistemas de control tolerante a fallas (FTC,
\textit{Fault-Tolerant Control}) para vehículos aéreos no tripulados (UAV) de tipo cuadricóptero ha adquirido gran relevancia en los ultimos años. En este contexto, garantizar una operación segura en entornos dinámicos representa un desafío. Por esta razón, es fundamental implementar estrategias de control capaces de detectar y compensar fallas, asegurando un funcionamiento estable y continuo del vehículo aun ante condiciones adversas. \cite{freddi2019,ozcan2024}.


\noindent Los UAV cuadricóptero se utilizan en aplicaciones de inspección, vigilancia, entrega, búsqueda y rescate, donde la pérdida de estabilidad ante una falla compromete tanto la misión como la seguridad del entorno en el que operan. Las consecuencias de una pérdida de control en sistemas multirotor que operan cerca de personas son reales: en diciembre de 2024, una pérdida de control simultánea en una formación de 500 drones durante un espectáculo en Orlando, Florida, provocó colisiones en el aire y la caída de varias unidades sobre el público, resultando en lesiones graves a un menor de edad, según el reporte preliminar de la Junta Nacional de Seguridad en el Transporte de Estados 
Unidos (NTSB) \cite{ntsb2024}. Si bien este incidente se originó por errores de configuración de la formación, evidencia la magnitud del riesgo que representa la pérdida de control en sistemas multirotor operando en proximidad a personas. La literatura especializada señala que, en escenarios operativos, las fallas de actuadores, sensores y estructura son frecuentes y pueden conducir a consecuencias igualmente desastrosas si no son detectadas y compensadas a tiempo \cite{quadcopter2024, fdd2021survey}. A diferencia de las aeronaves de ala fija, los cuadricópteros generan sustentación exclusivamente mediante la rotación de sus motores, lo que los hace 
vulnerables ante fallas de actuador: la pérdida de efectividad en un solo motor puede provocar una pérdida de control en cuestión de segundos \cite{amoozgar2013}.

\noindent Las fallas que afectan a estos sistemas se clasifican en 
tres categorías principales. Las fallas en actuadores corresponden 
a la pérdida de efectividad en uno o más motores, ya sea por 
desgaste mecánico, daño en hélices o sobrecalentamiento, y pueden 
representar pérdidas de empuje de hasta el 36\% en un motor 
individual sin llegar a una falla total \cite{kim2025, 
amoozgar2013}; en casos más severos, la pérdida completa de uno o 
varios rotores compromete directamente el estado del 
sistema, dado que un cuadricóptero tiene menos entradas de control
\cite{wang2022, du2023}. Las fallas en sensores incluyen lecturas 
erróneas o degradadas del giroscopio, el acelerómetro o el sistema 
de posicionamiento, las cuales pueden originarse por interferencia 
electromagnética, vibración mecánica excesiva o, en escenarios más 
críticos, por ataques ciberfísicos que falsifican deliberadamente 
las señales recibidas por la unidad de control 
\cite{patan2022, bekhiti2025}. A diferencia de las fallas de 
actuador, cuyo efecto es frecuentemente visible en la dinámica del 
vuelo, las fallas de sensor pueden pasar desapercibidas durante 
varios pasos de control, ya que el sistema continúa operando con 
información incorrecta hasta que la divergencia entre el estado 
estimado y el estado real se vuelve significativa. Finalmente, las 
perturbaciones externas de origen aerodinámico —como ráfagas de 
viento, efecto suelo o variaciones súbitas de carga— modifican la 
dinámica del sistema de forma impredecible y no pueden ser 
anticipadas mediante un modelo nominal, lo que exige mecanismos de 
estimación capaces de distinguir entre una perturbación externa 
transitoria y una falla persistente del propio sistema 
\cite{hsu2024}. Cada una de estas categorías afecta de manera 
distinta la dinámica del sistema y puede presentarse de forma 
aislada o combinada durante una misma trayectoria —por ejemplo, una 
ráfaga de viento puede coincidir con la pérdida parcial de un motor 
o con una lectura degradada del acelerómetro—, lo que exige un 
enfoque integrado capaz de atender los tres escenarios de forma 
simultánea.


\noindent Frente a esta problemática, los enfoques actuales abordan 
la detección y compensación de fallas mediante metodologías 
distintas, cuya complejidad radica en que cada una resuelve 
parcialmente el problema bajo supuestos particulares. Un primer 
grupo de trabajos emplea métodos basados en modelos matemáticos y 
observadores de estado, que generan señales residuales para 
detectar desviaciones respecto al comportamiento nominal del 
sistema \cite{freddi2019, patan2022, du2023}. Si bien estos métodos 
ofrecen eficiencia computacional, su efectividad depende 
directamente de la exactitud del modelo dinámico empleado, lo que 
limita su desempeño ante dinámicas no lineales o fallas simultáneas 
en múltiples componentes.


\noindent Como alternativa a los métodos basados en modelos, las 
técnicas de aprendizaje profundo permiten identificar patrones 
anómalos directamente a partir de señales de vuelo, sin requerir un 
modelo matemático explícito del sistema. Arquitecturas que combinan 
\textit{autoencoders} LSTM para reducción de dimensionalidad con 
clasificadores BiLSTM para análisis bidireccional de secuencias 
temporales han logrado predecir fallas hasta 30 segundos antes de 
que ocurra un evento crítico, con una precisión cercana al 80\% 
sobre datos de vuelo reales \cite{ozcan2024}. Desde una perspectiva 
semi-supervisada, modelos que combinan LSTM con \textit{One-Class 
SVM} (OCSVM) permiten detectar desviaciones anómalas en señales 
multivariadas de sensores sin necesidad de ejemplos etiquetados de 
falla durante el entrenamiento, lo cual resulta especialmente útil 
en escenarios donde los datos de falla son escasos \cite{li2021}. 
No obstante, ambos enfoques comparten una limitación estructural: 
se concentran en la tarea de detección sin incorporar un esquema de 
compensación activo que permita al UAV mantener la estabilidad una 
vez identificada la anomalía, y requieren grandes volúmenes de 
datos etiquetados o el ajuste cuidadoso de umbrales de decisión 
para evitar falsos positivos en condiciones operativas reales.

\noindent Por su parte, las políticas de control aprendidas mediante 
aprendizaje por refuerzo (RL), que han demostrado la 
capacidad de adaptarse a condiciones dinámicas no vistas durante el 
entrenamiento sin requerir un modelo explícito del sistema 
\cite{liu2024, sohege2021, kim2025}. Estos enfoques constituyen, en 
general, los de mejor desempeño reportado en la literatura; no 
obstante, la mayoría presenta brechas de transferencia 
simulación-realidad, y los que alcanzan los mejores resultados de 
adaptación se concentran en un único tipo de falla, sin incorporar 
un módulo explícito de detección y diagnóstico.
%


\noindent Esta brecha constituye una oportunidad de investigación 
concreta. Se ha probado el uso de arquitecturas con memoria 
temporal para abordar la detección y compensación de fallas en 
UAVs, principalmente mediante LSTM \cite{ozcan2024, li2021} y 
módulos basados en \textit{Transformer} \cite{kim2025, giral2024}; 
sin embargo, ambas presentan limitaciones para su despliegue en 
plataformas embebidas de bajo costo: las LSTM procesan la secuencia 
de forma estrictamente secuencial, lo que limita su 
paralelización, mientras que los Transformer requieren almacenar y 
atender sobre el historial completo de la secuencia, con un costo 
computacional y de memoria que crece cuadráticamente con la 
longitud de dicho historial. Por otro lado, la arquitectura Mamba, 
basada en modelos de espacio de estados selectivos, mantiene un 
estado oculto de tamaño constante durante la inferencia y procesa 
la secuencia mediante un mecanismo de \textit{scan} selectivo 
paralelizable, lo que reduce significativamente el costo 
computacional sin sacrificar la capacidad de modelar dependencias 
temporales largas. Podría ser especialmente relevante emplear Mamba 
en este contexto porque su perfil computacional resulta adecuado 
para la detección de fallas progresivas en UAVs, donde la anomalía 
se manifiesta gradualmente a lo largo de decenas de pasos 
temporales y donde las restricciones de hardware embebido limitan 
el tamaño de los modelos viables —condición bajo la cual ni las 
LSTM ni los Transformer han sido evaluados en la literatura 
revisada para esta tarea específica. Ante la necesidad identificada, 
se propone un sistema integrado de detección y compensación de 
fallas en UAV cuadricóptero mediante la combinación de estimación 
de estados y aprendizaje temporal ligero basado en Mamba, validado 
en el simulador gym-pybullet-drones con escenarios de falla en 
actuadores, sensores y perturbaciones externas, atendiendo así la 
brecha de integración y viabilidad computacional identificada en el 
estado del arte.




\introsection{Formulación del problema}

%Propuesta 1
De qué manera la integración de estimación de estados 
y aprendizaje temporal ligero basado en modelos de espacio de 
estados selectivos permite detectar y compensar fallas en sensores, 
actuadores y perturbaciones dinámicas en UAV cuadricópteros en 
tiempo real con viabilidad computacional para plataformas embebidas


%propuesta 2
¿Cómo un método inteligente y ligero de aprendizaje temporal permite detectar y compensar fallas de UAV cuadricópteros con viabilidad computacional?


\introsection{Objetivos de investigación}
Desarrollar un sistema integrado de 
detección y compensación de fallas en UAV cuadricóptero mediante 
estimación de estados y aprendizaje temporal ligero, validado en entorno simulado con escenarios de falla en 
actuadores, sensores y perturbaciones externas. \\

Los objetivos específicos son:
\begin{enumerate}
    \item  Diseñar un estimador de estados que detecte 
inconsistencias en las mediciones de sensores del UAV, evaluando 
su desempeño ante ruido y perturbaciones externas en el simulador 
gym-pybullet-drones.
    \item Desarrollar un modelo de detección de fallas basado en la arquitectura Mamba que capture patrones secuenciales y permita la identificación anticipada de anomalías en los tres escenarios definidos, comparando su desempeño frente a LSTM.

    \item  Implementar un esquema de compensación que mantenga la estabilidad del UAV ante fallas en actuadores y perturbaciones dinámicas, validando su desempeño en diferentes escenarios de operación simulados.
\end{enumerate}


\introsection{Justificación}
\noindent La presente investigación se justifica desde una 
perspectiva tecnológica y metodológica. Como se evidenció en la 
descripción de la situación problemática, los trabajos que alcanzan 
los mejores resultados de detección y compensación de fallas en UAV 
cuadricópteros emplean arquitecturas de aprendizaje profundo —LSTM y 
\textit{Transformer}— cuyo costo computacional dificulta su 
despliegue en plataformas embebidas de bajo costo, mientras que los 
trabajos validados en hardware embebido real se limitan a un único 
tipo de falla. Al integrar estimación de estados con una 
arquitectura de espacio de estados selectivo, este trabajo aporta un 
método de aprendizaje temporal cuyo costo computacional no crece con 
la longitud del historial procesado, lo que representa una 
alternativa metodológica frente a las arquitecturas predominantes en 
la literatura revisada. Asimismo, la validación conjunta de fallas 
en actuadores, sensores y perturbaciones externas dentro de un mismo 
marco aporta valor teórico al evidenciar si un enfoque integrado 
puede mantener un desempeño comparable al de los métodos 
especializados que abordan cada tipo de falla por separado.

\noindent En el plano práctico, los resultados de esta investigación 
podrían beneficiar a futuros trabajos de control tolerante a fallas 
en UAVs orientados a plataformas de bajo costo, al proporcionar una 
arquitectura cuyo perfil computacional la hace candidata para 
implementaciones embebidas, así como un conjunto de datos simulados 
con escenarios de falla en actuadores, sensores y perturbaciones 
aerodinámicas, generado en gym-pybullet-drones, que puede ser 
reutilizado en investigaciones posteriores sobre detección y 
compensación de fallas en cuadricópteros. De manera más amplia, dado 
que los UAV cuadricóptero se emplean en aplicaciones de inspección, 
vigilancia y entrega que operan en proximidad a personas e 
infraestructura, contar con métodos de detección y compensación de 
fallas computacionalmente viables contribuye a reducir el riesgo de 
incidentes asociados a la pérdida de control en este tipo de 
sistemas.


\introsection{Alcance y limitaciones / restricciones} %[opcional]

El \textbf{alcance}  comprende el desarrollo 
y validación de un sistema integrado de detección y compensación 
de fallas en un UAV cuadricóptero simulado mediante 
gym-pybullet-drones \cite{panerati2021}, utilizando el modelo 
Crazyflie CF2X como plataforma de referencia, con tres escenarios 
de falla: pérdida de efectividad en actuadores (motores), 
corrupción de señales de sensores (giroscopio) y perturbaciones 
externas de viento lateral. La validación se realiza 
exclusivamente en entorno simulado, reconociendo la brecha 
simulación-realidad (\textit{sim-to-real gap}) como trabajo 
futuro \cite{kim2025}, \cite{hsu2024}; la física de PyBullet 
constituye una aproximación de la dinámica real y no modela 
efectos aerodinámicos complejos como turbulencias ni interacciones 
entre rotores \cite{panerati2021}; y la evaluación de viabilidad 
en hardware embebido se limita al análisis del costo computacional 
del modelo entrenado, sin despliegue físico.